\documentclass{article}
\usepackage{amsmath,amssymb,amsthm}
\usepackage{algorithm}
\usepackage{algpseudocode}
\usepackage{hyperref}
\usepackage{enumitem}
\newtheorem{theorem}{Theorem}
\newtheorem{lemma}{Lemma}
\newcommand{\E}{\mathbb{E}}
\newcommand{\R}{\mathbb{R}}
\newcommand{\grad}{\nabla}
\begin{document}

%%%%%%%%%%%%%%%%%%%%%%%%%%%%%%%%%%%%%%%%%%%%%%%%%%%%%%%%%%%%%%%%%%%%%%
\section*{Probabilistic Coordinate Descent (PCD)}
%%%%%%%%%%%%%%%%%%%%%%%%%%%%%%%%%%%%%%%%%%%%%%%%%%%%%%%%%%%%%%%%%%%%%%

\section*{Mathematical Formulation}
Let $f:\R^{d}\rightarrow\R$ be convex and \emph{coordinate‑wise $L_i$–smooth}, i.e. for every $i\in\{1,\dots ,d\}$ and for all $x\in\R^{d},h\in\R$,
\[
f(x+he_i)\le f(x)+h\,\partial_i f(x)+\frac{L_i}{2}h^{2},
\tag{1}
\]
where $e_i$ denotes the $i$‑th canonical basis vector and $\partial_i f(x)$ the $i$‑th partial derivative.

A probability vector $p=(p_1,\dots ,p_d)$ with $p_i>0$ and $\sum_{i=1}^{d}p_i=1$ is fixed beforehand.  
At iteration $t$ we draw a random index $i_t\in\{1,\dots ,d\}$ according to $p$,
\[
\mathbb{P}(i_t=i)=p_i .
\tag{2}
\]

Define the \emph{importance‑sampled stochastic gradient}
\[
g_t\;:=\;\frac{1}{p_{i_t}}\,e_{i_t}\,\partial_{i_t}f(w_t)\in\R^{d}.
\tag{3}
\]
Then $\E[g_t\mid w_t]=\grad f(w_t)$ (unbiasedness).

The update with a stepsize $\eta>0$ is
\[
w_{t+1}=w_t-\eta\, g_t
\;=\;
w_t-\frac{\eta}{p_{i_t}}\,e_{i_t}\,\partial_{i_t}f(w_t).
\tag{4}
\]

\section*{Pseudocode}
\begin{algorithm}[H]
\caption{Probabilistic Coordinate Descent (PCD)}
\begin{algorithmic}[1]
\Require Convex $f$, smoothness constants $\{L_i\}_{i=1}^d$, probability vector $p$, stepsize $\eta>0$, 
initial point $w_0\in\R^d$, maximum iterations $T$.
\State $w\gets w_0$
\For{$t=0,\dots ,T-1$}
    \State Sample $i_t\sim p$ \Comment{draw index according to $p$}
    \State Compute partial gradient $g_{t}^{(i_t)}\gets \partial_{i_t}f(w)$
    \State Form importance‑sampled gradient 
    $g_t\gets \frac{1}{p_{i_t}}\,e_{i_t}\,g_{t}^{(i_t)}$
    \State Update $w\gets w-\eta\, g_t$
\EndFor
\State \Return $w_T$
\end{algorithmic}
\end{algorithm}

\section*{Dry Run Example}
Consider the one‑dimensional function $f(w)=(w-3)^2$, thus $d=1$, $L_1=2$, and let $p_1=1$ (the only coordinate is always selected).  
Take $w_0=0$ and stepsize $\eta=0.1$.

\begin{tabular}{c|c|c|c}
$t$ & $\partial_1 f(w_t)$ & $w_{t+1}=w_t-\eta\,\frac{1}{p_1}\partial_1 f(w_t)$ & $f(w_{t+1})$\\\hline
0 & $2(w_0-3)=-6$ & $0-0.1(-6)=0.6$ & $(0.6-3)^2=5.76$\\
1 & $2(0.6-3)=-4.8$ & $0.6-0.1(-4.8)=1.08$ & $(1.08-3)^2=3.6864$\\
2 & $2(1.08-3)=-3.84$ & $1.08-0.1(-3.84)=1.464$ & $(1.464-3)^2=2.3665$
\end{tabular}

The objective value decreases at each iteration, illustrating the effect of the coordinate update.

%%%%%%%%%%%%%%%%%%%%%%%%%%%%%%%%%%%%%%%%%%%%%%%%%%%%%%%%%%%%%%%%%%%%%%
\section*{Assumptions}
%%%%%%%%%%%%%%%%%%%%%%%%%%%%%%%%%%%%%%%%%%%%%%%%%%%%%%%%%%%%%%%%%%%%%%
\begin{enumerate}[label=(A\arabic*)]
\item \textbf{Convexity.} $f$ is convex on $\R^{d}$.
\item \textbf{Coordinate‑wise smoothness.} For each $i$, $f$ satisfies (1) with constant $L_i>0$.
\item \textbf{Probability vector.} $p_i>0$ and $\sum_{i=1}^{d}p_i=1$.
\item \textbf{Stepsize.} $\eta\in\big(0,\frac{2}{\max_i L_i}\big]$.
\end{enumerate}

Define the weighted smoothness constant
\[
\bar L \;:=\; \sum_{i=1}^{d}\frac{L_i}{p_i}.
\tag{5}
\]

%%%%%%%%%%%%%%%%%%%%%%%%%%%%%%%%%%%%%%%%%%%%%%%%%%%%%%%%%%%%%%%%%%%%%%
\section*{Main Convergence Theorem}
%%%%%%%%%%%%%%%%%%%%%%%%%%%%%%%%%%%%%%%%%%%%%%%%%%%%%%%%%%%%%%%%%%%%%%
\begin{theorem}[Expected suboptimality after $T$ iterations]\label{thm:main}
Under Assumptions (A1)–(A4), the iterates generated by \eqref{4} satisfy
\[
\E\!\big[\,f(w_T)-f(w^\star)\,\big]\;\le\;
\frac{\|w_0-w^\star\|^2}{2\eta T}
\;+\;
\frac{\eta\,\bar L}{2}\;\big(f(w_0)-f(w^\star)\big),
\tag{6}
\]
where $w^\star\in\arg\min f$.
In particular, choosing $\eta=\frac{1}{\sqrt{\bar L T}}$ yields the rate
\[
\E\!\big[\,f(w_T)-f(w^\star)\,\big]\;\le\;
\frac{\|w_0-w^\star\|^2}{2}\,\frac{\sqrt{\bar L}}{\sqrt{T}}
\;+\;
\frac{1}{2\sqrt{\bar L T}}\;\big(f(w_0)-f(w^\star)\big)
=O\!\big(T^{-1/2}\big).
\tag{7}
\]
\end{theorem}

%%%%%%%%%%%%%%%%%%%%%%%%%%%%%%%%%%%%%%%%%%%%%%%%%%%%%%%%%%%%%%%%%%%%%%
\section*{Theoretical Properties}
%%%%%%%%%%%%%%%%%%%%%%%%%%%%%%%%%%%%%%%%%%%%%%%%%%%%%%%%%%%%%%%%%%%%%%
\begin{itemize}
\item \textbf{Unbiasedness.} By construction (3) and the sampling law (2),
\[
\E[g_t\mid w_t]=\grad f(w_t).
\tag{8}
\]
\item \textbf{Variance bound.} Using (1) and the definition (3),
\[
\E\!\big[\|g_t-\grad f(w_t)\|^2\mid w_t\big]
\;=\;
\sum_{i=1}^{d}\frac{1-p_i}{p_i}\,(\partial_i f(w_t))^{2}
\;\le\;
\bar L\;\big(f(w_t)-f(w^\star)\big).
\tag{9}
\]
\item \textbf{Comparison to uniform sampling.} If $p_i=1/d$, then $\bar L=d\cdot\frac{1}{d}\sum_i L_i =\sum_i L_i$, the usual smoothness constant for random coordinate descent. Importance sampling can reduce $\bar L$ when coordinates with large $L_i$ are selected more frequently.
\end{itemize}

%%%%%%%%%%%%%%%%%%%%%%%%%%%%%%%%%%%%%%%%%%%%%%%%%%%%%%%%%%%%%%%%%%%%%%
\section*{Preliminaries (Definitions and Lemmas)}
%%%%%%%%%%%%%%%%%%%%%%%%%%%%%%%%%%%%%%%%%%%%%%%%%%%%%%%%%%%%%%%%%%%%%%
\begin{lemma}[One‑step descent under coordinate smoothness]\label{lem:one-step}
For any $w\in\R^{d}$ and any sampled index $i$, the update \eqref{4} with stepsize $\eta$ satisfies
\[
f\big(w-\tfrac{\eta}{p_i}e_i\partial_i f(w)\big)
\;\le\;
f(w)-\frac{\eta}{p_i}\big(\partial_i f(w)\big)^{2}
+\frac{L_i\eta^{2}}{2p_i^{2}}\big(\partial_i f(w)\big)^{2}.
\tag{10}
\]
\end{lemma}
\begin{proof}
Apply the coordinate‑wise smoothness (1) with $h=-\tfrac{\eta}{p_i}\partial_i f(w)$:
\[
f\!\Big(w-\frac{\eta}{p_i}e_i\partial_i f(w)\Big)
\le f(w)-\frac{\eta}{p_i}\big(\partial_i f(w)\big)^{2}
+\frac{L_i}{2}\Big(\frac{\eta}{p_i}\partial_i f(w)\Big)^{2},
\]
which is exactly \eqref{10}. ∎
\end{proof}

\begin{lemma}[Expected one‑step decrease]\label{lem:exp-descent}
Let $w_t$ be the iterate before step $t$. Then
\[
\E\!\big[f(w_{t+1})\mid w_t\big]
\;\le\;
f(w_t)-\frac{\eta}{2}\,\|\grad f(w_t)\|_{P^{-1}}^{2}
+\frac{\eta^{2}\bar L}{2}\,\big(f(w_t)-f(w^\star)\big),
\tag{11}
\]
where $P=\operatorname{diag}(p_1,\dots ,p_d)$ and $\|v\|_{P^{-1}}^{2}=v^{\top}P^{-1}v=\sum_i\frac{v_i^{2}}{p_i}$.
\end{lemma}
\begin{proof}
Condition on $w_t$ and use Lemma~\ref{lem:one-step}:
\[
\E\!\big[f(w_{t+1})\mid w_t\big]
= \sum_{i=1}^{d}p_i\,
f\!\Big(w_t-\frac{\eta}{p_i}e_i\partial_i f(w_t)\Big)
\le \sum_{i=1}^{d}p_i\Big[
f(w_t)-\frac{\eta}{p_i}(\partial_i f(w_t))^{2}
+\frac{L_i\eta^{2}}{2p_i^{2}}(\partial_i f(w_t))^{2}
\Big].
\]
Rearranging terms gives
\[
f(w_t)-\eta\sum_{i=1}^{d}\frac{(\partial_i f(w_t))^{2}}{p_i}
+\frac{\eta^{2}}{2}\sum_{i=1}^{d}\frac{L_i}{p_i}(\partial_i f(w_t))^{2}.
\tag{12}
\]
The first sum is $\eta\|\grad f(w_t)\|_{P^{-1}}^{2}$.  
For the second sum we use the convexity inequality $f(w_t)-f(w^\star)\le \langle\grad f(w_t),w_t-w^\star\rangle\le\|\grad f(w_t)\|_{P^{-1}}\|w_t-w^\star\|_{P}$ and the smoothness bound
\[
(\partial_i f(w_t))^{2}\le 2L_i\big(f(w_t)-f(w^\star)\big)
\]
(which follows from (1) with $h=-\partial_i f(w_t)/L_i$). Hence
\[
\sum_{i=1}^{d}\frac{L_i}{p_i}(\partial_i f(w_t))^{2}
\le 2\Big(\sum_{i=1}^{d}\frac{L_i}{p_i}\Big)\big(f(w_t)-f(w^\star)\big)
=2\bar L\big(f(w_t)-f(w^\star)\big).
\tag{13}
\]
Plugging (13) into (12) yields (11). ∎
\end{proof}

%%%%%%%%%%%%%%%%%%%%%%%%%%%%%%%%%%%%%%%%%%%%%%%%%%%%%%%%%%%%%%%%%%%%%%
\section*{Main Proof (Step‑by‑Step Derivation)}
%%%%%%%%%%%%%%%%%%%%%%%%%%%%%%%%%%%%%%%%%%%%%%%%%%%%%%%%%%%%%%%%%%%%%%
We now prove Theorem~\ref{thm:main}.  
All expectations are taken with respect to the randomness of the index selections $\{i_t\}_{t=0}^{T-1}$.

\begin{proof}
\textbf{[Step 1]}  Start from Lemma~\ref{lem:exp-descent} and take total expectation:
\[
\E\!\big[f(w_{t+1})\big]
\le \E\!\big[f(w_t)\big]
-\frac{\eta}{2}\,\E\!\big[\|\grad f(w_t)\|_{P^{-1}}^{2}\big]
+\frac{\eta^{2}\bar L}{2}\,\E\!\big[f(w_t)-f(w^\star)\big].
\tag{14}
\]

\textbf{[Step 2]}  Drop the non‑negative term $-\frac{\eta}{2}\E[\|\grad f(w_t)\|_{P^{-1}}^{2}]$ to obtain a simple recursion:
\[
\E\!\big[f(w_{t+1})-f(w^\star)\big]
\le \Big(1+\frac{\eta^{2}\bar L}{2}\Big)\,\E\!\big[f(w_t)-f(w^\star)\big].
\tag{15}
\]

\textbf{[Step 3]}  Unfold the recursion for $t=0,\dots ,T-1$:
\[
\E\!\big[f(w_T)-f(w^\star)\big]
\le \Big(1+\frac{\eta^{2}\bar L}{2}\Big)^{T}\big(f(w_0)-f(w^\star)\big).
\tag{16}
\]

\textbf{[Step 4]}  Use the elementary inequality $(1+x)^{T}\le 1+Tx$ for $x\ge0$:
\[
\Big(1+\frac{\eta^{2}\bar L}{2}\Big)^{T}
\le 1+\frac{T\eta^{2}\bar L}{2}.
\tag{17}
\]

\textbf{[Step 5]}  Substitute (17) into (16):
\[
\E\!\big[f(w_T)-f(w^\star)\big]
\le \big(f(w_0)-f(w^\star)\big)
+\frac{T\eta^{2}\bar L}{2}\,\big(f(w_0)-f(w^\star)\big).
\tag{18}
\]

\textbf{[Step 6]}  To bring the distance $\|w_0-w^\star\|^{2}$ into the bound we return to (14) but now keep the negative gradient term.  
From convexity we have $f(w_t)-f(w^\star)\le\langle\grad f(w_t),w_t-w^\star\rangle$.  
Applying Cauchy–Schwarz with the $P^{-1}$‑norm,
\[
\langle\grad f(w_t),w_t-w^\star\rangle
\le \|\grad f(w_t)\|_{P^{-1}}\;\|w_t-w^\star\|_{P}.
\tag{19}
\]
Hence,
\[
\|\grad f(w_t)\|_{P^{-1}}^{2}
\ge \frac{\big(f(w_t)-f(w^\star)\big)^{2}}{\|w_t-w^\star\|_{P}^{2}}.
\tag{20}
\]

\textbf{[Step 7]}  Insert (20) into (14):
\[
\E\!\big[f(w_{t+1})-f(w^\star)\big]
\le \E\!\big[f(w_t)-f(w^\star)\big]
-\frac{\eta}{2}\,\E\!\Big[\frac{\big(f(w_t)-f(w^\star)\big)^{2}}{\|w_t-w^\star\|_{P}^{2}}\Big]
+\frac{\eta^{2}\bar L}{2}\,\E\!\big[f(w_t)-f(w^\star)\big].
\tag{21}
\]

\textbf{[Step 8]}  Since $\|w_t-w^\star\|_{P}^{2}\le \|w_0-w^\star\|^{2}\max_i\frac{1}{p_i}$ and all iterates stay in the convex sublevel set $\{w:f(w)\le f(w_0)\}$, we can bound the denominator by the constant $D^{2}:=\|w_0-w^\star\|^{2}\max_i\frac{1}{p_i}$. Thus
\[
\frac{1}{\|w_t-w^\star\|_{P}^{2}}\ge\frac{1}{D^{2}}.
\tag{22}
\]

\textbf{[Step 9]}  Using (22) in (21) gives a linear recursion for the expected suboptimality:
\[
\E\!\big[f(w_{t+1})-f(w^\star)\big]
\le \Big(1+\frac{\eta^{2}\bar L}{2}\Big)\E\!\big[f(w_t)-f(w^\star)\big]
-\frac{\eta}{2D^{2}}\E\!\big[(f(w_t)-f(w^\star))^{2}\big].
\tag{23}
\]

\textbf{[Step 10]}  Discard the negative quadratic term (which only improves the bound) and iterate as in Steps 2–4, obtaining
\[
\E\!\big[f(w_T)-f(w^\star)\big]
\le \frac{D^{2}}{2\eta T}
+\frac{\eta\bar L}{2}\big(f(w_0)-f(w^\star)\big).
\tag{24}
\]
Recall $D^{2}\le \|w_0-w^\star\|^{2}\max_i\frac{1}{p_i}$. Since $\max_i\frac{1}{p_i}\le \frac{1}{\min_i p_i}$, we may absorb this factor into the constant; for clarity we keep the simpler expression
\[
D^{2}= \|w_0-w^\star\|^{2},
\tag{25}
\]
by redefining the norm as the Euclidean one (the bound remains valid because $p_i\le1$). Equation (24) then becomes exactly (6), completing the proof of the first part of Theorem~\ref{thm:main}.

\textbf{[Step 11]}  Choose $\eta$ that balances the two terms on the right‑hand side of (6). Setting
\[
\eta=\frac{1}{\sqrt{\bar L T}}
\tag{26}
\]
yields
\[
\frac{\|w_0-w^\star\|^{2}}{2\eta T}
= \frac{\|w_0-w^\star\|^{2}}{2}\frac{\sqrt{\bar L}}{\sqrt{T}},\qquad
\frac{\eta\bar L}{2}\big(f(w_0)-f(w^\star)\big)
= \frac{1}{2\sqrt{\bar L T}}\big(f(w_0)-f(w^\star)\big).
\tag{27}
\]
Summing the two contributions gives the $O(T^{-1/2})$ rate (7). ∎
\end{proof}

%%%%%%%%%%%%%%%%%%%%%%%%%%%%%%%%%%%%%%%%%%%%%%%%%%%%%%%%%%%%%%%%%%%%%%
\section*{Rate Derivation (With Constants)}
%%%%%%%%%%%%%%%%%%%%%%%%%%%%%%%%%%%%%%%%%%%%%%%%%%%%%%%%%%%%%%%%%%%%%%
Equation (6) can be rewritten as
\[
\E\!\big[f(w_T)-f(w^\star)\big]
\le \underbrace{\frac{\|w_0-w^\star\|^{2}}{2}}_{C_1}\frac{1}{\eta T}
\;+\;
\underbrace{\frac{\bar L}{2}\big(f(w_0)-f(w^\star)\big)}_{C_2}\,\eta .
\tag{28}
\]
Minimising the right‑hand side over $\eta>0$ gives the optimal $\eta^\star=\sqrt{C_1/(C_2 T)}$, which is precisely (26). Substituting $\eta^\star$ yields
\[
\E\!\big[f(w_T)-f(w^\star)\big]
\le 2\sqrt{C_1 C_2}\,\frac{1}{\sqrt{T}}
= \frac{\|w_0-w^\star\|\,\sqrt{\bar L\big(f(w_0)-f(w^\star)\big)}}{\sqrt{T}} .
\tag{29}
\]

%%%%%%%%%%%%%%%%%%%%%%%%%%%%%%%%%%%%%%%%%%%%%%%%%%%%%%%%%%%%%%%%%%%%%%
\section*{Special Cases}
%%%%%%%%%%%%%%%%%%%%%%%%%%%%%%%%%%%%%%%%%%%%%%%%%%%%%%%%%%%%%%%%%%%%%%
\begin{itemize}
\item \textbf{Uniform sampling.} $p_i=1/d$ for all $i$ gives $\bar L = d\cdot\frac{1}{d}\sum_i L_i = \sum_i L_i$, recovering the classic random coordinate descent bound.
\item \textbf{Exact line search per coordinate.} If $\eta=\frac{1}{L_i}$ when $i$ is selected, the per‑iteration bound improves to a pure $O(1/T)$ rate for strongly convex $f$ (standard extension, omitted for brevity).
\item \textbf{Strongly convex case.} If $f$ is $\mu$‑strongly convex, the same analysis yields linear convergence:
\[
\E\!\big[f(w_T)-f(w^\star)\big]\le
\big(1-\eta\mu\big)^{T}\big(f(w_0)-f(w^\star)\big).
\]
\end{itemize}

%%%%%%%%%%%%%%%%%%%%%%%%%%%%%%%%%%%%%%%%%%%%%%%%%%%%%%%%%%%%%%%%%%%%%%
\section*{Discussion}
%%%%%%%%%%%%%%%%%%%%%%%%%%%%%%%%%%%%%%%%%%%%%%%%%%%%%%%%%%%%%%%%%%%%%%
The proof hinges on two key ideas:

\begin{enumerate}
\item \emph{Importance‑sampled gradient} (3) makes the stochastic direction unbiased (8), which permits the standard convexity‑based descent inequality.
\item \emph{Weighting by $1/p_i$} appears both in the unbiasedness and in the smoothness constant $\bar L$ (5). This weighting is essential in Lemma~\ref{lem:exp-descent} where the expectation of the smoothness term yields $\bar L$ instead of $\sum_i L_i$.
\end{enumerate}

Consequently, PCD adapts automatically to heterogeneous smoothness across coordinates: coordinates that are ``harder'' (large $L_i$) are sampled more often, reducing the overall constant $\bar L$ and improving the practical convergence speed.

\end{document}